\begin{abstract}
The Goal-Frontier Maximization framework's central measure
$\volL$ quantifies capability \emph{possession}, but the
experiential value it protects requires capability
\emph{exercise}~($\volR$).  Global direct measurement of $\volR$
would require panopticon-scale observation, which the framework's
privacy commitments rule out.  We introduce
\emph{revealed-sacrifice observation}: under assumption S1
(free choice with non-degrading alternatives), each voluntary
trade event---an agent surrendering a benchmarked $X$ for an
unbenchmarked bundle $Y$---reveals the preference comparison
$\Ui(Y) \geq \Ui(X)$.  Combined with S0
(valuation-bounded-by-exercise; an assumed bridge from
agent valuation to realized exercise, not itself observed),
S2 (benchmark grounding), and S4(a) (valuation additivity),
the preference comparison extends to a realized-contribution
lower bound $\volR^{[W]}(Y) \geq \Delta\volL(X)$.  The bound
is a theorem about the pairing of observation and
assumption-stack, not a direct disclosure of $\volR$ content.  Aggregating across events with
pairwise poset-independent bundles yields the \emph{Aggregate
B-to-C Lower Bound theorem}: on the observed subset $U$ within a
trade window $W$,
$\volR^{U,[W]} \geq \sum_n \Delta\volL(X_n)$,
under assumptions S0--S2, S4(a), and the genuine-exchange
conditions of Definition~\ref{def:sacrifice_event}.  The bound
is non-decreasing in observed events, satisfies three formal
privacy properties, and composes with the commitment protocol
of~\citep{lasser2026exo} to produce two interfaces: a default
distinct-category batch interface that publishes per-event
$(\Delta\volL, \mathrm{category})$ pairs, and a hiding-variant
rollup that certifies only an aggregate threshold over a
hidden subset.  Both interfaces' privacy profiles depend on
two named institutional preconditions
(Proposition~\ref{prop:privacy_institutional}) in addition to
the formal properties.  Money and time
sacrifice channels triangulate the same underlying substrate.
We also show that the mapping
$\mathcal{P} \mapsto \volR^{\mathcal{P}}$ from capability
systems to the realized volume is \emph{not strictly monotone
under capability removal}: removing an unexercised capability
leaves $\volR$ unchanged while $\volL$ strictly contracts.  The
self-balancing property of~\citep[Theorem~1]{lasser2026poset}
therefore does not transfer to $\volR$ (this is a system-level
mapping property, not a measure-axiom failure on a fixed
exercised sub-poset, where $\volR$ inherits M1--M6 by
restriction).  Consequently $\volR$ is a diagnostic, not an
optimization target.  The need-sufficiency filter enabling S1 across events,
the gap decomposition, wireheading-consistent concentration
signal, and residual class are developed
in~\citep{lasser2026paper8b}.
\end{abstract}
